\section{Related Work}
\label{sec:related}

\paragraph{Shadow detection in overhead imagery.}
Overhead shadow detection assigns a binary shadow label to every pixel of an aerial or satellite image. The field was established by AISD~\cite{aisd}, a 514-image manually annotated benchmark spanning five cities at 0.3\,m resolution. S-EO~\cite{seo} scales to roughly 19{,}000 images by deriving masks from digital surface models and sun geometry; the authors document temporal-misalignment artifacts and recommend the dataset for training rather than evaluation. Airborne-Shadow~\cite{azimi2023airborne} adds a manually annotated overhead set with a heavy/light shadow distinction, but is restricted to German imagery and does not report a geographic split. Recent methods (MAMNet, MBCSANet, CSDFormer, OGLANet) report F1 scores within roughly one point of each other on AISD~\cite{mamnet, mbcsanet, csdformer, oglanet}, indicating that the in-domain benchmark has plateaued. Additionally, ground-level shadow datasets~\cite{istd, sbu, cuhkshadow} are not substitutes for overhead evaluation: viewing angle, and shadow geometry differ. None of these overhead datasets defines a geographic train-test split, and the most recent overhead shadow survey identifies cross-dataset generalization as an open problem~\cite{hu2026unveiling}.

\paragraph{Domain adaptation for dense prediction.}
A range of methods make dense prediction robust to source-target appearance differences by intervening at one stage of the pipeline. At the input, FDA swaps low-frequency Fourier amplitude between images~\cite{fda}, IIM derives illumination-invariant inputs~\cite{yola}, and BSM mixes styles within a training batch through a frozen AdaIN network~\cite{luo2023bsm}. In feature space, ISW whitens photometrically sensitive covariance dimensions~\cite{isw}, SegDeSiC aligns encoder features through a geo-coordinate regression loss~\cite{segdesicnet}, MRFP$^+$ randomizes high- and low-frequency feature components~\cite{mrfp}, and FADA processes the Haar-wavelet bands of foundation-model features separately~\cite{fada}. At test time, TENT minimises prediction entropy over normalisation-layer parameters~\cite{wang2021tent}, and SAR stabilises this update with sharpness-aware minimisation and by excluding high-entropy samples~\cite{niu2023sar}. Each method applies its correction at the same locus for every architecture, and no prior work tests whether that locus is consistent across architectures for overhead dense prediction. Remote-sensing foundation models address the shift through the representation instead: CrossEarth trains adapters and a Mask2Former decoder on a frozen DINOv2 backbone, with style augmentation from an external image pool and masked image modelling~\cite{gong2025crossearth}.

